\section{Introduction}
\label{sec:introduction}

India's Business Responsibility and Sustainability Reporting framework gives environmental, social, and governance disclosures a common structure. A 2021 circular from the Securities and Exchange Board of India (SEBI) made BRSR reporting mandatory for the top 1,000 listed companies by market capitalisation from financial year 2022--23. The framework covers nine principles of responsible business conduct. Essential indicators are mandatory, while leadership indicators are voluntary \cite{sebi2021}. A common format helps organise reports, but cannot establish whether every response is complete or correct.

A data analyst must complete several manual steps when reviewing a company report. A text segment is a selected portion of report text used for analysis. First, a relevant text segment must be found and assigned to a reporting topic. Next, a relevant numerical value must be located and checked against the reporting year, unit, and business activities covered. Comparing companies adds another challenge because business activities and applicable reporting items differ. Repeating these steps across hundreds of reports takes substantial time. For example, a statement that \emph{water use decreased} needs water figures for the relevant years and the same business operations before the claim can be assessed.

Automated text analysis can reduce this workload. ClimateBERT identifies climate-related language and selected disclosure categories. ChatReport retrieves relevant text segments from reports and answers reporting questions with source links. ESGReveal extracts structured information from narrative reports \cite{webersinke2021,ni2023,zou2025}. These methods help organise large report collections. A complete review workflow must also explain how claims are checked, how training labels are obtained, and how company scores are calculated \cite{calamai2025,garg2025}. These requirements lead to three questions.

\textbf{First, how can a written claim be connected with supporting evidence?} Recognising that a text segment discusses water use does not establish whether a claim of reduced water use is supported. The relevant reported values must also be available for review. GRASP stores text segments and reported values under the same company report. Language checks mark vague wording for review. Numerical checks produce records containing the relevant sentence and reported values. Assurance checks add information about external verification reported by the company. These records help an analyst examine the evidence behind a flag.

\textbf{Second, how can training labels be obtained from the reports?} A classifier learns from text segments paired with labels. Each label identifies one of the nine BRSR principles. GRASP obtains labels from three sources: explicit principle tags, mappings between reporting questions and principles, and numbered principle notes. Some question mappings are AI-assisted. Labels derived through these rules are called \emph{weak labels}. GRASP records the source of each label and compares classifier results across the three sources. This record explains how the training data are constructed and what the reported classification results measure.

\textbf{Third, what does a company's Transparency Score represent?} A total score becomes easier to understand when the calculation and score components are visible. GRASP reports four components. Completeness measures the coverage of reporting items. Specificity measures the share of selected text segments containing evidence markers, such as numbers or recognised standards. Clarity measures the share without a vagueness flag. The fourth component counts Sustainable Development Goals connected to text segments with evidence markers. These connections follow predefined links between BRSR principles and the goals. GRASP shows the contribution of each component to the total score and compares companies within the same recorded sector. An analyst can therefore examine both the score breakdown and the position among sector peers.

The study uses 42,409 labelled text segments. Companies are separated into training, validation, and test groups. Training teaches the classifier to assign principles. Validation helps select a model version during development. Testing measures performance on companies excluded from supervised training. Additional tests use text segments from reporting items excluded from classifier training. These reporting items are called \emph{concepts} in the structured data. This comparison examines performance on reporting questions outside the training set.

The language detector is tested with settings kept unchanged during evaluation. Numerical checks create review records that show a claim alongside the relevant reported figures. The Transparency Score is compared with previously collected human ratings. Results are presented step by step: one company is examined first, followed by selected financial companies, the full financial-sector group, and all recorded sectors. Each evaluation explains what is being tested and which labels or ratings are used for comparison. This workflow brings report text, numerical evidence, assigned BRSR principles, and score calculations together to help an analyst review sustainability disclosures.

The remaining sections are organised as follows. Section~\ref{sec:literature} reviews related research. Section~\ref{sec:proposed} presents the proposed solution. Section~\ref{sec:methodology} explains how each of the seven modules works. Section~\ref{sec:evaluation} describes the evaluation data and model results. Section~\ref{sec:results} presents company and sector comparisons. Section~\ref{sec:conclusion} summarises the findings and future work.
